\documentclass[10pt,a4paper]{amsart}
\usepackage[margin=19mm]{geometry}
\usepackage{amsmath,amssymb,mathtools}
\usepackage[scr=rsfs,cal=euler]{mathalfa}
\usepackage{xfrac}
\usepackage[unicode,hidelinks,bookmarks=false]{hyperref}
\newcommand{\ie}{i.e.\ }
\newcommand{\cf}{cf.\ }
\newcommand{\resp}{respectively\ }
\newcommand{\Subsection}{\subsection}
\providecommand{\nobreakdash}{\nobreak\hbox{-}}
\setlength{\emergencystretch}{2em}
\multlinegap0pt
\usepackage{mathtools}
\usepackage{amssymb}
\usepackage{float}
\usepackage{algorithm}\usepackage{algpseudocode}
\usepackage[normalem]{ulem}
\usepackage{xcolor}
\newtheorem{lemma}{Lemma}
\def\Rbkk{{\mathbb{R}^{k\times k}}}
\def\Rbnk{{\mathbb{R}^{n\times k}}}
\def\Rbnn{{\mathbb{R}^{n\times n}}}
\def\iStkn{\mathrm{iSt}(k,n)}
\title{A Riemannian gradient descent method\\for optimization on the indefinite Stiefel manifold}
\author{Dinh Van Tiep, Nguyen Thanh Son}
\date{Source: arXiv:2410.22068v3 (CC BY 4.0)}
\begin{document}
\maketitle

\noindent\emph{Source and licence.} A Riemannian gradient descent method\\for optimization on the indefinite Stiefel manifold. Version arXiv:2410.22068v3 (CC BY 4.0), distributed by arXiv under the Creative Commons Attribution 4.0 International licence, http://creativecommons.org/licenses/by/4.0/, as recorded in the arXiv licence metadata for this exact version and shown on the abstract page https://arxiv.org/abs/2410.22068v3. Full licence notice: \texttt{licenses/arxiv-2410.22068-CC-BY-4.0.txt}.\par\smallskip

\noindent\textit{Editorial scope.} Counted unit: the article's lemma prop:eleproperty (source0.tex lines 544--567). The article's complete original proof of it is delivered with it (source0.tex lines 568--644, one \texttt{proof} environment of 77 lines). No mathematics is written, repaired or re-derived: the statement and the proof are the article's own lines, in the article's own notation, and no retained line is substituted or re-flowed.  No further source-local context is needed: every cross-reference of every retained line resolves inside the two delivered spans.  Nothing was omitted from any retained span, so the package carries no omission record.  No retained line contains a citation, so the wrapper carries no bibliography and no bibliography entry.  No retained line contains a figure environment, an image-inclusion command or drawing code, so no image is delivered and the package declares no asset dependency.  Editorial formatting only, in addition: an AMS \texttt{amsart} preamble that restates the article's own declarations and macros used by the retained lines, namely the environments \texttt{lemma} on separate counters, together with the standard packages those lines need.  The retained environments print in this document as Lemma 1 in order of appearance, whereas the article prints the same items with the numbering of its own full document.  The complete native source of the article as distributed in the arXiv e-print for this exact version is bundled unchanged as \texttt{sources/167593/source0.tex}.
			\begin{lemma}\label{prop:eleproperty} 
				For $X\in \iStkn$, let us denote the matrix
				\begin{equation*}%\label{eq:Ematrix}
					E\coloneqq [X \;\; A^{-1}X_{\perp}]\in \Rbnn.
				\end{equation*} Then
				\begin{itemize}
					\item[i)] $E$ is nonsingular; 
					\item[ii)] $E^{T}AE = 
					\begin{bmatrix}
						J & 0 \\ 
						0 & X_{\perp}^{T}A^{-1}X_{\perp} 
					\end{bmatrix}$ and $X_{\perp}^{T}A^{-1}X_{\perp}$ is nonsingular;
					\item[iii)] $E^{-1}=
					\begin{bmatrix}
						JX^{T}A \\ 
						(X_{\perp}^{T}A^{-1}X_{\perp})^{-1}X_{\perp}^{T} 
					\end{bmatrix};$
					\item[iv)] For any $Z\in\Rbnk,$ there are $W\in\Rbkk$ and $K\in\mathbb{R}^{(n-k)\times k} $ such that
					\begin{equation}
						Z = XW + A^{-1}X_{\perp}K.\label{eq:eq21}
					\end{equation}
					Moreover, $W = JX^{T}AZ, K = (X_{\perp}^{T}A^{-1}X_{\perp})^{-1}X_{\perp}^{T}Z,$ and $A^{-1}X_{\perp}K$ is independent of $X_{\perp}.$
				\end{itemize} 	
			\end{lemma}

			\begin{proof}
				\begin{itemize}\item[i)]
					To proceed, we will verify that the homogeneous linear system $E \left[\begin{smallmatrix}
							y_1 \\ 
							y_2
						\end{smallmatrix}\right]=0$ %\in\mathbb{R}^{n\times 1},$ 
						with $y_1\in\mathbb{R}^{k}, y_2\in\mathbb{R}^{n-k}$ has only the trivial solution. %This can be seen when
					Multiply both sides of the equation by $X^{T}A$ from the left to obtain %$X^{T}AX_1 + X^{T}AA^{-1}X_{\perp}Y_2 = 0$ or 
						$Jy_1 = 0$. This results in %since $X^{T}X_{\perp} = 0.$ Hence, $Y_1 =0.$ This follows that the homogeneous system is just 
					$A^{-1}X_{\perp}y_2 = 0$ which immediately implies that $y_2=0$ %This means $Y_2 = 0$ 
					because $A^{-1}X_{\perp}\in\mathbb{R}^{n\times (n-k)}$ is of full rank.% matrix.
					\item[ii)] The identity follows from a direct calculation. %$E^{T}AE = E^{T} [AX\;X_{\perp}] = \begin{bmatrix}
						%			X^{T} \\ 
						%			X_{\perp}^{T}A^{-T}
						%		\end{bmatrix} [AX\;X_{\perp}] = \begin{bmatrix}
						%		X^{T}AX & X^{T}X_{\perp} \\ 
						%		X_{\perp}^{T}A^{-1}AX & X_{\perp}^{T}A^{-1}X_{\perp}
						%		\end{bmatrix} = \begin{bmatrix}
						%		J & 0\\ 
						%		0 & X_{\perp}^{T}A^{-1}X_{\perp}
						%		\end{bmatrix}.$ 
					Moreover, since the matrices $E$ and $A$ are invertible, so are $E^{T}AE$ and its diagonal block $X_{\perp}^{T}A^{-1}X_{\perp}.$
					\item[iii)] The equality in ii) yields that 
						$	E^{-1} = \left[\begin{smallmatrix}
							J & 0\\ 
							0 & (X_{\perp}^{T}A^{-1}X_{\perp})^{-1}
						\end{smallmatrix}\right] E^TA
						$
						which leads to the expected equality. %From ii, multiplying both sides of that identity to the right by $E^{-1}$ and to the left by the inverse matrix of the right hand side of the identity yields
					%		$  \begin{bmatrix}
						%			J & 0\\ 
						%			0 & (X_{\perp}^{T}A^{-1}X_{\perp})^{-1}
						%		\end{bmatrix}E^{T}A = E^{-1}$ or $ \begin{bmatrix}
						%		J & 0\\ 
						%		0 & (X_{\perp}^{T}A^{-1}X_{\perp})^{-1}
						%		\end{bmatrix} \begin{bmatrix}
						%		X^{T} A\\ 
						%		 X_{\perp}^{T}
						%		\end{bmatrix} = E^{-1}.$
					%		The last identity means that
					%		$E^{-1} = \begin{bmatrix}
						%			JX^{T} A\\ 
						%			(X_{\perp}^{T}A^{-1}X_{\perp})^{-1}X_{\perp}^{T}
						%		\end{bmatrix}.$
					\item[iv)] For any $Z\in\Rbnk,$ let $W\in\Rbkk$ and $K\in\mathbb{R}^{(n-k)\times k}$ be the matrices obtained from the identity $E^{-1}Z =  \left[\begin{smallmatrix}
						W\\ 
						K
					\end{smallmatrix}\right].$ From iii), we get $W = JX^{T}AZ$ and $ K = (X_{\perp}^{T}A^{-1}X_{\perp})^{-1}X_{\perp}^{T}Z.$ %By definition, we have
					So,
					$Z = E  \left[\begin{smallmatrix}
						W\\ 
						K
					\end{smallmatrix}\right] = [X\; A^{-1}X_{\perp}] \left[\begin{smallmatrix}
						W\\ 
						K
					\end{smallmatrix}\right] = XW +  A^{-1}X_{\perp}K.$ 
					The independence of $X_{\perp}$  %$A^{-1}X_{\perp}K$ from  
						is derived from the fact %facts 
						that $A^{-1}X_{\perp}K = Z - XW$ with $W = JX^TAZ$.
				\end{itemize}
				% is independent from $X_{\perp}$.
				%		Moreover, from the identity 
				%		\begin{equation}\label{eq:eq22}
					%		\begin{aligned}
						%			I_{n} = E^{-1}E &= [X\; \;A^{-1}X_{\perp}]\begin{bmatrix}
							%				JX^{T} A\\ 
							%				(X_{\perp}^{T}A^{-1}X_{\perp})^{-1}X_{\perp}^{T}
							%			\end{bmatrix} \\
						%			&= XJX^{T}A + A^{-1}X_{\perp}(X_{\perp}^{T}A^{-1}X_{\perp})^{-1}X_{\perp}^{T},
						%		\end{aligned}
					%		\end{equation}
				%        it follows that
				%		\begin{center}
					%		$A^{-1}X_{\perp}K = A^{-1}X_{\perp}(X_{\perp}^{T}A^{-1}X_{\perp})^{-1}X_{\perp}^{T}Z = (I_n - XJX^{T}A)Z.$
					%		\end{center}
				%		Therefore, $A^{-1}X_{\perp}K$ is independent of $X_{\perp}.$
			\end{proof}

\end{document}
